\documentclass[11pt,a4paper]{article}
\usepackage[T1]{fontenc}
\usepackage[utf8]{inputenc}
\usepackage[margin=22mm]{geometry}
\usepackage{lmodern}
\usepackage{amsmath,amssymb}
\usepackage{enumitem}
\usepackage{xcolor}
\usepackage{microtype}
\usepackage{hyperref}
\definecolor{epflred}{HTML}{E3000B}
\hypersetup{colorlinks=true,urlcolor=epflred,linkcolor=epflred,
  pdftitle={ENG-654 Project 18: Local versus Global Joint-Limit Avoidance on iiwa 7}}
\setlength{\parindent}{0pt}
\setlength{\parskip}{6pt}
\setlist[itemize]{leftmargin=*,itemsep=5pt,topsep=4pt}

\begin{document}
{\small\textbf{EPFL \textbar\ ENG-654}\hfill Kinematics-Grounded Motion Planning for Robots}
\vspace{5mm}

{\color{epflred}\Large\bfseries Project 18}

{\LARGE\bfseries Local versus Global Joint-Limit Avoidance on iiwa 7\par}

\textbf{Format:} Group of two students; simulation-based project.\\
\textbf{Prerequisites:} Redundant inverse kinematics, null-space projection, joint-limit avoidance, and feasibility-map planning.

\section*{Motivation}
Null-space motion can improve a local joint-centering objective while
missing a feasible route elsewhere in the redundant IK set. Comparing that
local strategy with a global feasibility chart helps distinguish algorithmic
limitations from actual restrictions imposed by the robot and the task.

\section*{Task}
\textbf{Overall objectives.} Compare local null-space joint centering with planning through a
classified iiwa 7 IK chart, and explain success or failure through feasible
self-motion intervals and their connections along the path.

\begin{itemize}
  \item \textbf{Establish one common problem.} Validate the iiwa model, limits,
and complete pose path. Use one short path segment that approaches a joint
limit and two admissible initial configurations, ensuring both methods are
compared from the same initial configuration in each trial.

  \item \textbf{Implement local joint centering.} Use the geometric Jacobian with
a documented scaling and a projected gradient for
$h(q)=\tfrac12\sum_i((q_i-q_{i,c})/r_i)^2$.
Integrate small steps, correct the full-pose error, and recheck limits.
Explain why quadratic centering is a preference, not a hard limit barrier.

  \item \textbf{Build the global comparison chart.} Enumerate the supplied
fixed-$q_3$ IK families on path-sample versus redundant-angle axes. Preserve
branch labels, plot feasible intervals and failure reasons, and distinguish
full-Jacobian rank loss from reduced-chart degeneracy.

  \item \textbf{Plan under the same constraints.} Select a chart corridor from the
same starting configuration, using joint centering and joint travel as
documented costs. Validate intermediate transitions and apply the same
pose, limit, and regularity tolerances used for the local method.

  \item \textbf{Diagnose the difference.} Overlay both trajectories on the
classification charts. Compare completion, minimum limit margin, travel,
and pose error for the two initial states. Repeat the local method with
one smaller integration step to distinguish discretization error from
failure to discover another feasible corridor.
\end{itemize}

\newpage
\section*{Specifics and scope}
Use one path segment, two initial configurations, and two local step
sizes. Start with about 120 path samples and 120 redundancy angles and
refine decisive intervals. Reuse the supplied analytical IK and construct the chart by
enumerating and validating candidates on the path--redundancy grid; a complete global optimizer is unnecessary. Equal results
are acceptable when the explanation is supported by the feasibility maps.

For every chart, keep pose validity, joint limits,
full-Jacobian regularity, and reduced-chart regularity as separate diagnostics.
The rectangular $6\times7$ Jacobian has no determinant. A feasible pixel is
not a validated edge: re-solve intermediate prescribed poses and check joint
continuity within limits. If collision checking is unavailable, state that
feasibility is kinematic only; a clearance proxy is not a full collision test.

\section*{Expected outcome}
Understand the difference between instantaneous null-space freedom
and global feasible connectivity. Deliver classified redundancy charts with
trajectory overlays, matched comparisons, and an explanation of whether
observed limitations arise from the strategy, numerical steps, or constraints.

Submit reproducible code, model parameters and test data, the required plots,
a concise 5--6-page report, and a short simulation demonstration. Explain
both successful cases and failures; conclusions must follow the numerical
and geometric evidence rather than an expected outcome.

\section*{Resources}
The KUKA LBR iiwa 7 robot description (URDF) and visualization meshes
(STL) will be provided to students, together with a mathematical continuation
model with unrestricted joints and an analytical fixed-$q_3$ IK model. Reuse
the supplied IK rather than deriving a general 7R solver. Apply the physical
joint limits from the robot description separately when using the unrestricted
model. Check D--H parameters, joint axes, and the fixed tool transform against
the URDF, and validate IK outputs by independent forward kinematics.

A sampled rectangular full-pose path and a reference fixed-$q_3$ IK
table will also be provided. The pose data contain position coordinates and
quaternion columns \texttt{qw,qx,qy,qz}; normalize the quaternions and make
their signs consistent before interpolation. The reference IK table checks
solutions at $q_3=30^\circ$ only; it does not replace evaluation at other angles.
Use its configurations to check pose residuals and branch labels, applying
physical joint limits separately.

Implement the null-space controller and feasibility-chart comparison using
the supplied IK and model. Choose a short segment or documented placement of
the supplied path; a suitable joint-limit example can be supplied to keep
scenario construction bounded. Visualize both recovered trajectories with
the supplied meshes and use identical validation tolerances.

\end{document}
